\section{About the course}

\subsection{History of the course}
\begin{frame}
\frametitle{History of the course}
This course was \alert{newly} developed by \textbf{Anthony R. Thornton} for the 2025/2026 academic year.

It is based on:
\begin{description}
\item[Mathematics 2M1 | Probability and Statistics] written by Robert Gaunt, School of Natural Sciences, University of Manchester,
 and David J. Silvester, School of Natural Sciences, University of Manchester.
\end{description}

\only<presentation>{\vspace{0.3cm}}
\textit{Special thanks to Google Gemini for proofreading and structural suggestions.}
\end{frame}

\subsection{Additional reading}
\begin{frame}[allowframebreaks]
\frametitle{Optional Books}
\only<article>{It should be noted that this reader, with associated lecture slides forms the \emph{entire} content of the course and is complete. However for the interested student a selection of text books is includes some following the same style as the course and other that are more mathematically rigourous. It should be noted all the resources listed below are full \textbf{open access}.}

\only<presentation>{Optional \textbf{Open Access} additional reading for this course:}

\begin{block}{Primary Resources}
\begin{itemize}
    \item \textbf{Module 1 (Descriptive Statistics):} \\
    \emph{Introductory Statistics} by \textbf{OpenStax}. \\
    \url{https://openstax.org/details/books/introductory-statistics}

    \item \textbf{Modules 2--4 (Probability):} \\
    \emph{Introduction to Probability, Statistics, and Random Processes} \\
    by \textbf{H. Pishro-Nik}. (Tailored for EEE). \\
    \url{https://www.probabilitycourse.com}
\end{itemize}
\end{block}

\begin{block}{Supplementary / General Reference}
\begin{itemize}
    \item \emph{OpenIntro Statistics} by \textbf{Diez, Çetinkaya-Rundel, \& Barr}. \\
    (A clean, concise alternative for general statistics concepts). \\
    \url{https://www.openintro.org/book/os/}

    \item \emph{Think Stats} by \textbf{Allen B. Downey}. \\
    (Teaches statistics through Python programming). \\
    \url{https://greenteapress.com/wp/think-stats-2e/}

    \item \textbf{Khan Academy (Statistics \& Probability)}. \\
    (Video tutorials useful for visualising histograms and distributions). \\
    \url{https://www.khanacademy.org/math/statistics-probability}
\end{itemize}
\end{block}

\begin{block}{For the Mathematically Interested}
    \emph{Introduction to Probability} by \textbf{Grinstead \& Snell}. \\
    (Classic text with rigorous proofs and derivations). \\
    \url{https://math.dartmouth.edu/~prob/prob/prob.pdf}
\end{block}
\end{frame}

\subsection{Course Structure}
\begin{frame}
\frametitle{Course Structure}
This course is designed to cover the final part of Mathematics for EEE 1E2, namely:

\begin{description}
    \item[Statistics] Discrete and continuous data. Mean, mode, and median. Range, variance and \textbf{standard deviation}. Frequency curves: Normal distribution and standardised normal.
    \item[\& Probability] Empirical and classical. Addition and multiplication laws. Probability distributions. Mean and standard deviation. Binomial and Poisson distribution. Normal distribution curves.
\end{description}

It will consist of four approximately 90-minute lectures.
\only<article>{The course is written in a highly practical way, focusing on the application of the topic rather than proofs, derivations, and strict mathematical rigour. However, it also teaches you how to identify the underlying mathematics in real-world scenarios, using many examples drawn from electrical engineering. Ultimately, this is not a pure maths course; it is a maths \textit{for} engineering course.}

\only<presentation>{\vspace{0.5cm} % Adds a little breathing room
}

The first two elements: Differential Vector Calculus and Integral Vector Calculus are \textbf{covered} by a separate course taught by \alert{Igor Chernyavsky}.
\end{frame}

\subsection{Assessment}
\begin{frame}
\frametitle{Assessment}
The assessment for these elements of the course consists of:
\begin{itemize}
    \setlength\itemsep{1em}
    \item \textbf{Online Coursework} (Week 10) \\
    \hspace{0.5cm} Contributes \textbf{10\%} to the final grade.
    
    \item \textbf{End-of-Semester Examination} \\
    \hspace{0.5cm} Contributes \textbf{80\%} to the final grade.
\end{itemize}

\vspace{0.5cm} % Adds separation for the final note
\small{The remaining \textbf{10\%} is assessed in the section taught by \textbf{Igor Chernyavsky}.}
\only<presentation>{\\[2mm]}
Note both these test wills also include questions from the earlier parts 1 and 2.
\end{frame}

\subsection{Course Overview}
\begin{frame}
\frametitle{Course Overview}
    \textbf{Overview:} 5 Thematic Blocks delivered across 4 Lectures (100 mins each).
    
    \vspace{0.3cm}
    \begin{center}
    \small
    \begin{tabular}{|l|l|l|}
        \hline
        \textbf{Session} & \textbf{Thematic Blocks} & \textbf{Learning Outcomes} \\ \hline
        Lecture 1 & Overview of Part &  \\
         & \textbf{Block 1:} Descriptive Statistics & ILO 5: Data Analysis \\ \hline
        Lecture 2 & \textbf{Block 2:} Prob. Fundamentals & ILO 6: Modeling \\ 
                  & \textbf{Block 3:} Discrete Distributions & \\ \hline
        Lecture 3 & \textbf{Block 4:} Continuous Distributions & ILO 6: Problem Solving \\ \hline
        Lecture 4 & \textbf{Block 5:} Linear Regression & ILO 5: Interpretation \\
                   & \\ \hline
    \end{tabular}
    \end{center}

    
\end{frame}

\begin{frame}[fragile]{Accessing the Living Course Materials}
    While official PDFs are on \textbf{Canvas}, the \lq Living\rq\ version of this course (including live Python code and LaTeX source) is hosted on GitHub.
    
    \vspace{0.3cm}
    \begin{block}{GitHub Repository}
        \centering
        \url{https://github.com/antRthorn/UoM-MATH19622-Public}
    \end{block}

    \vspace{0.3cm}
    \structure{How to get the code (For the 10\%):}
    \begin{enumerate}
        \item \textbf{Clone the Repo:}
        \begin{lstlisting}[language=bash]
git clone https://github.com/antRthorn/UoM-MATH19622-Public.git
        \end{lstlisting}
        \item \textbf{Run the Examples:}
        \begin{lstlisting}[language=bash]
uv run path/to/script.py
        \end{lstlisting}
    \end{enumerate}
    \footnotesize \textit{Note: Use the VS Code Dev Container for a zero-install Python environment.}
\end{frame}

\only<presentation>{
\begin{frame}{The QA Bounty Program}
      \begin{itemize}
        \item \textbf{The Task:} Find a typo, a calculation error, or suggest a major improvement to the course code or examples.
        \item \textbf{The Rewards:}
            \begin{itemize}
                \item \textbf{\pounds 5 Voucher:} For minor errors (typos, broken links).
                \item \textbf{\pounds 10 Voucher:} For significant clarity improvements or bug fixes.
            \end{itemize}
        \item \textbf{The Process:} Submit a \textbf{Pull Request} or an \textbf{Issue} on GitHub.
        \item \textbf{The Catch:} First person to report the specific error wins the bounty for that slide!
    \end{itemize}
    \vspace{0.3cm}
    \centering
    \textit{\lq An engineer who never makes a mistake has never made anything.\rq}
\end{frame}
}

